\input{preamble}

\section*{Mott scattering}

Mott scattering is the result of kinetic electrons $e^-$ interacting with atomic nuclei $Z$.

\begin{center}
\begin{tikzpicture}
\draw[dashed] (0,0) circle (0.5cm);
\draw (0,0) node {$Z$};
\draw[thick,->] (-2,0) node[anchor=east] {$e^-$} -- (-0.6,0);
\draw[thick,->] (0.40,0.40) -- (1.3,1.3) node[anchor=south west] {$e^-$};
\draw (1,0.5) node {$\theta$};
\end{tikzpicture}
\end{center}

Momentum vectors where $p$ is electron momentum and $E$ is relativistic energy $E=\sqrt{p^2+m^2}$.
\begin{equation*}
p_1=(E,\mathbf p_1)
=\underset{\substack{\text{inbound}\\ \text{electron}}}
{\begin{pmatrix}E\\0\\0\\p\end{pmatrix}}
\qquad
p_2=(E,\mathbf p_2)
=\underset{\substack{\text{outbound}\\ \text{electron}}}
{\begin{pmatrix}
E\\
p\sin\theta\cos\phi\\
p\sin\theta\sin\phi\\
p\cos\theta
\end{pmatrix}}
\end{equation*}

Spinors for the inbound electron. All spinors are normalized as $u^\dag u=2E$.
\begin{equation*}
u_{11}=\frac{1}{\sqrt{E+m}}
\underset{\substack{\text{inbound electron}\\ \text{spin up}}}
{\begin{pmatrix}E+m\\0\\p\\0\end{pmatrix}}
\qquad
u_{12}=\frac{1}{\sqrt{E+m}}
\underset{\substack{\text{inbound electron}\\ \text{spin down}}}
{\begin{pmatrix}0\\E+m\\0\\-p\end{pmatrix}}
\end{equation*}

Spinors for the outbound electron.
\begin{equation*}
u_{21}=\frac{1}{\sqrt{E+m}}
\underset{\substack{\text{outbound electron}\\ \text{spin up}}}
{\begin{pmatrix}E+m\\0\\p_{2z}\\p_{2x}+ip_{2y}\end{pmatrix}}
\qquad
u_{22}=\frac{1}{\sqrt{E+m}}
\underset{\substack{\text{outbound electron}\\ \text{spin down}}}
{\begin{pmatrix}0\\E+m\\p_{2x}-ip_{2y}\\-p_{2z}\end{pmatrix}}
\end{equation*}

Inbound wave functions.
All wave functions are normalized as $\langle\psi|\psi\rangle=1$ over volume $V$.
\begin{align*}
\psi_{11}(\mathbf r)&=\frac{1}{\sqrt{2EV}}u_{11}\exp\left(\frac{i}{\hbar}\mathbf p_1\cdot\mathbf r\right)
\\[1ex]
\psi_{12}(\mathbf r)&=\frac{1}{\sqrt{2EV}}u_{12}\exp\left(\frac{i}{\hbar}\mathbf p_1\cdot\mathbf r\right)
\end{align*}

Outbound wave functions.
\begin{align*}
\psi_{21}(\mathbf r)&=\frac{1}{\sqrt{2EV}}u_{21}\exp\left(\frac{i}{\hbar}\mathbf p_2\cdot\mathbf r\right)
\\[1ex]
\psi_{22}(\mathbf r)&=\frac{1}{\sqrt{2EV}}u_{22}\exp\left(\frac{i}{\hbar}\mathbf p_2\cdot\mathbf r\right)
\end{align*}

For the Coulomb potential
\begin{equation*}
V(\mathbf r)=-\frac{Ze^2}{4\pi\epsilon_0r}=-\frac{Z\alpha}{r}
\end{equation*}

we have the following transition amplitudes.
\begin{align*}
M_1&=-Z\alpha
\int\frac{1}{r}\psi_{21}^\dag(\mathbf r)\psi_{11}(\mathbf r)\,d^3\mathbf r
& &\text{up $\leftarrow$ up}
\\[1ex]
M_2&=-Z\alpha
\int\frac{1}{r}\psi_{21}^\dag(\mathbf r)\psi_{12}(\mathbf r)\,d^3\mathbf r
& &\text{up $\leftarrow$ down}
\\[1ex]
M_3&=-Z\alpha
\int\frac{1}{r}\psi_{22}^\dag(\mathbf r)\psi_{11}(\mathbf r)\,d^3\mathbf r
& &\text{down $\leftarrow$ up}
\\[1ex]
M_4&=-Z\alpha
\int\frac{1}{r}\psi_{22}^\dag(\mathbf r)\psi_{12}(\mathbf r)\,d^3\mathbf r
& &\text{down $\leftarrow$ down}
\end{align*}

Substitute for the wave functions.
\begin{align*}
M_1&=-\frac{1}{2EV}u_{21}^\dag u_{11}Z\alpha
\int\frac{1}{r}
\exp\left[\frac{i}{\hbar}(\mathbf p_1-\mathbf p_2)\cdot\mathbf r\right]\,d^3\mathbf r
\\[1ex]
M_2&=-\frac{1}{2EV}u_{21}^\dag u_{12}Z\alpha
\int\frac{1}{r}
\exp\left[\frac{i}{\hbar}(\mathbf p_1-\mathbf p_2)\cdot\mathbf r\right]\,d^3\mathbf r
\\[1ex]
M_3&=-\frac{1}{2EV}u_{22}^\dag u_{11}Z\alpha
\int\frac{1}{r}
\exp\left[\frac{i}{\hbar}(\mathbf p_1-\mathbf p_2)\cdot\mathbf r\right]\,d^3\mathbf r
\\[1ex]
M_4&=-\frac{1}{2EV}u_{22}^\dag u_{12}Z\alpha
\int\frac{1}{r}
\exp\left[\frac{i}{\hbar}(\mathbf p_1-\mathbf p_2)\cdot\mathbf r\right]\,d^3\mathbf r
\end{align*}

Solve the integral to obtain
\begin{equation*}
\int\frac{1}{r}
\exp\left[\frac{i}{\hbar}(\mathbf p_1-\mathbf p_2)\cdot\mathbf r\right]\,d^3\mathbf r
=\frac{\pi\hbar^2}{p^2\sin^2(\theta/2)}
\tag{1}
\end{equation*}

Hence
\begin{align*}
M_1&=-\frac{1}{2EV}u_{21}^\dag u_{11}
\frac{Z\alpha\pi\hbar^2}{p^2\sin^2(\theta/2)}
\\[1ex]
M_2&=-\frac{1}{2EV}u_{21}^\dag u_{12}
\frac{Z\alpha\pi\hbar^2}{p^2\sin^2(\theta/2)}
\\[1ex]
M_3&=-\frac{1}{2EV}u_{22}^\dag u_{11}
\frac{Z\alpha\pi\hbar^2}{p^2\sin^2(\theta/2)}
\\[1ex]
M_4&=-\frac{1}{2EV}u_{22}^\dag u_{12}
\frac{Z\alpha\pi\hbar^2}{p^2\sin^2(\theta/2)}
\end{align*}

Expected density $\langle|M|^2\rangle$ is the average density.
(There are two inbound states hence the divisor is 2.)
\begin{equation*}
\langle|M|^2\rangle=\frac{1}{2}\left(|M_1|^2+|M_2|^2+|M_3|^2+|M_4|^2\right)
\end{equation*}

Factor the densities.
\begin{equation*}
\langle|M|^2\rangle
=\left(\frac{1}{2EV}
\frac{Z\alpha\pi\hbar^2}{p^2\sin^2(\theta/2)}\right)^2
\frac{1}{2}
\left(|u_{21}^\dag u_{11}|^2+|u_{21}^\dag u_{12}|^2+|u_{22}^\dag u_{11}|^2+|u_{22}^\dag u_{12}|^2\right)
\end{equation*}

By computer algebra
\begin{equation*}
\frac{1}{2}
\left(|u_{21}^\dag u_{11}|^2+|u_{21}^\dag u_{12}|^2+|u_{22}^\dag u_{11}|^2+|u_{22}^\dag u_{12}|^2\right)
=4m^2+4p^2\cos^2(\theta/2)
\tag{2}
\end{equation*}

Hence
\begin{equation*}
\langle|M|^2\rangle
=\left(\frac{1}{2EV}
\frac{Z\alpha\pi\hbar^2}{p^2\sin^2(\theta/2)}\right)^2
\left(4m^2+4p^2\cos^2(\theta/2)\right)
\end{equation*}

By Fermi's golden rule
\begin{equation*}
\frac{d\sigma}{d\Omega}=\left(\frac{EV}{2\pi\hbar^2}\right)^2\langle|M|^2\rangle
\end{equation*}

Hence
\begin{equation*}
\frac{d\sigma}{d\Omega}
=\left(\frac{Z\alpha}{4p^2\sin^2(\theta/2)}\right)^2
\left(4m^2+4p^2\cos^2(\theta/2)\right)
\tag{3}
\end{equation*}

Factor the $4m^2$ to obtain
\begin{equation*}
\frac{d\sigma}{d\Omega}
=\left(\frac{Z\alpha m}{2p^2\sin^2(\theta/2)}\right)^2
\left(1+v^2\cos^2(\theta/2)\right)
\tag{4}
\end{equation*}

where $v=p/m$.
Noting that
\begin{equation*}
\underset{\text{Rutherford}}{\frac{d\sigma}{d\Omega}}
=\left(\frac{Z\alpha m}{2p^2\sin^2(\theta/2)}\right)^2
\end{equation*}

we have
\begin{equation*}
\frac{d\sigma}{d\Omega}
=\underset{\text{Rutherford}}{\frac{d\sigma}{d\Omega}}
\left(1+v^2\cos^2(\theta/2)\right)
\end{equation*}

\href{https://georgeweigt.github.io/examples/mott-scattering-demo.html}{Eigenmath script}

\end{document}
